% Blue River Dam — Implementation Surface
% PhD Thesis Chapter: Lifecycle Authority and Governance Engine

\section{Implementation Surface}
\label{sec:blue-river-dam-implementation}

\subsection{Source Files}
\label{sec:blue-river-dam-implementation-sources}

The implementation surface is intentionally minimal. The entire library is
contained in a single source file:

\begin{description}
  \item[\texttt{blue\_river\_dam/src/lib.rs}] 965 lines as of 2026-06-01
    (post-receipt state). Organized into five sections delimited by comment
    banners: (1) Governance Authority Hierarchy, (2) MAPE-K Loop
    Implementation, (3) Lifecycle State Machine, (4) Orchestrator root struct
    with \texttt{mape\_k\_cycle()}, and (5) Test suite. The file carries the
    module-level docstring citing the three primary governance doctrine sources
    from which it was synthesized, and the declaration
    \texttt{License: Executable only under wasm4pm graduation bridge}.
  \item[\texttt{blue\_river\_dam/Cargo.toml}] Minimal Cargo manifest: package
    name \texttt{blue\_river\_dam}, version 0.1.0, edition 2021. Empty
    \texttt{[dependencies]} and \texttt{[dev-dependencies]} sections, confirming
    zero external dependencies.
  \item[\texttt{blue\_river\_dam/README.md}] Project-level documentation
    (content not cited in receipt).
  \item[\texttt{blue\_river\_dam/GENERATION\_RECEIPT.md}] The render
    certification document. Records 5/5 tests passing, zero unsafe code, zero
    warnings, and full MAPE-K loop closure verification.
\end{description}

A root-level artifact, \texttt{liblib.rlib} (95{,}336 bytes, dated
2026-06-01T00:09), also exists in the project directory. This appears to be a
pre-Cargo compilation artifact from a direct \texttt{rustc} invocation
(as documented in \texttt{GENERATION\_RECEIPT.md}'s compilation trace), predating
the Cargo-based build. Its relationship to the Cargo-built rlibs is
undocumented; see Section~\ref{sec:blue-river-dam-questions-gaps}.

\subsection{Commands}
\label{sec:blue-river-dam-implementation-commands}

The project follows the global \texttt{cargo make} preference documented in
\texttt{$\sim$/.claude/CLAUDE.md}. However, because the \texttt{Cargo.toml}
contains no \texttt{[workspace]} declaration and no \texttt{Makefile.toml},
the following standard Cargo commands are the authoritative build surface:

\begin{itemize}
  \item \texttt{cargo build --release} --- produces
    \texttt{target/release/libblue\_river\_dam.rlib} and
    \texttt{target/release/libblue\_river\_dam.dylib} (cdylib target).
  \item \texttt{cargo test --lib} --- compiles and runs all 8 test functions
    in the \texttt{\#[cfg(test)] mod tests} block.
  \item \texttt{cargo check} --- type-checks without linking; the fastest
    way to verify that no \texttt{unsafe} block has been introduced.
  \item \texttt{rustc --crate-type lib src/lib.rs --edition 2021 -W unsafe-code}
    --- the direct-invocation form used during initial synthesis, as recorded
    in \texttt{GENERATION\_RECEIPT.md}.
\end{itemize}

\subsection{Data and Ontology Surfaces}
\label{sec:blue-river-dam-implementation-data}

Blue River Dam does not expose an RDF or SPARQL surface directly. Its data
model is expressed entirely in Rust types. The key data-bearing types are:

\begin{itemize}
  \item \texttt{PetriNetTopology\{places: Vec<String>, transitions: Vec<String>,
    flow: Vec<(String, String)>\}} --- the input type to WF-net soundness
    validation.
  \item \texttt{EventStream\{events: Vec<ProcessEvent>, window\_size: usize\}} ---
    the runtime event input type consumed by \texttt{Monitor::ingest\_stream()}.
  \item \texttt{ProcessEvent\{timestamp: u64, activity: u32, case\_id: u64\}} ---
    the atomic unit of process evidence.
  \item \texttt{ConformanceMetric\{fitness: f64, trace\_id: u64, alignment\_moves: u32,
    threshold: f64\}} --- the primary conformance output type.
  \item \texttt{Knowledge\{reference\_model: String, historical\_metrics: String,
    violation\_patterns: String, successful\_repairs: String,
    repair\_success\_count: u32, repair\_failure\_count: u32\}} --- the
    persistent MAPE-K knowledge store.
\end{itemize}

The autonomic activation layer extends this surface with YAML-format governance
artifacts: \texttt{AndonGatePolicy} (5 gates), \texttt{DriftDetectionMonitor}
(6 SPARQL-based rules), and a \texttt{MAPEKLoop} manifest --- all recorded in
\texttt{AUTONOMIC\_ACTIVATION\_LOG\_20260601.yaml}.

\subsection{Templates and Rendered Artifacts}
\label{sec:blue-river-dam-implementation-templates}

Blue River Dam is itself a rendered artifact, not the renderer. The
\texttt{GENERATION\_RECEIPT.md} documents that the library was synthesized
by the Lifecycle Orchestrator Renderer consuming five governance doctrine
sources. The rendering strategy is documented in the receipt under four
headings: Authority Hierarchy, MAPE-K Loop, Lifecycle State Machine, and
Governance Enforcement.

The 11 Jinja2 templates that make downstream rendering possible are housed in
\texttt{wasm4pm-compat/templates/} rather than in this project directory.
The lifecycle state machine template
(\texttt{templates/lifecycle/state\_machine.rs.j2}) and the actuation template
(\texttt{templates/lifecycle/actuation.rs.j2}) are the upstream templates
whose logical structure is instantiated by the Blue River Dam library.
[AUTHOR THESIS: the formal correspondence between the Jinja2 templates and the
synthesized Rust source has not been mechanically verified in this project
directory; it is asserted by the render receipt.]
